\section{Results Available at This Stage}
\label{sec:results}

This section reports the two blocks that are complete. Both were obtained
with \emph{simulator-oracle} obstacle perception rather than the visual
detector, and both are labelled accordingly throughout: they validate the
simulated plant and isolate the temporal estimator, and they are not visual
obstacle-avoidance results. The planner comparison (block C), the calibration
ladder (block D), and the physical campaign (block E) are not reported here.

\subsection{Simulation Integration Baseline}

The integration campaign accepted 9 of 9 runs across its three scenarios
under pre-registered criteria (Table~\ref{tab:baseline0}). Every avoidance
run produced the same state sequence with exactly one maneuver and no side
reversal, no collision occurred, minimum clearance to the obstacle surface
stayed in 0.672--0.715\,m, the vehicle returned to within 6\,cm of its
original line, and the transferable dead-reckoning error stayed at or below
0.20\,m. In the dropout scenario a deterministic 2\,s silence injected
mid-maneuver was survived in all three runs with no premature recovery, no
side switch, and no collision. Repeatability across fresh processes is tight
rather than exact --- 4\,cm spread in clearance, 0.3\,s in maneuver time ---
which is the expected consequence of distributed wall-clock scheduling.

%==============================================================================
% Scientific Baseline 0.  Source: docs/SCIENTIFIC_BASELINE_0.md,
% experiments/simulation/scientific_baseline_0/ (manifest + per-run JSON).
% Status: READY.  Label: simulation dynamics integration baseline, ORACLE
% perception -- NOT a visual obstacle-avoidance result.
%==============================================================================
\begin{table}[t]
\centering
\caption{Simulation integration baseline, 9/9 accepted runs from fresh
processes. Obstacle perception is the simulator oracle: this validates
physics, control, navigation and safety plumbing, \emph{not} vision.
Clearance is to the obstacle surface, computed from known geometry.}
\label{tab:baseline0}
\footnotesize
\setlength{\tabcolsep}{3pt}
\renewcommand{\arraystretch}{1.15}
\begin{tabular}{@{}p{26mm} r r r@{}}
\toprule
\textbf{Metric} & \textbf{A: straight} & \textbf{B: avoid} &
\textbf{C: avoid$+$dropout} \\
 & ($n{=}3$) & ($n{=}3$) & ($n{=}3$) \\
\midrule
Accepted                    & 3/3 & 3/3 & 3/3 \\
Collisions                  & --- & 0/3 & 0/3 \\
Min.\ clearance [m]         & --- & 0.672--0.712 & 0.688--0.715 \\
Forward progress [m]        & 13.4--14.0 & 34.9--35.0 & 34.4--35.0 \\
Max.\ lateral dev.\ [m]     & 0.000 & 2.43--2.48 & 2.45--2.47 \\
Final lateral error [m]     & 0.000 & 0.019--0.058 & 0.016--0.045 \\
Returned to line            & --- & 3/3 & 3/3 \\
Avoidance entries / switches& 0 / 0 & 1 / 0 & 1 / 0 \\
Maneuver time [s]           & --- & 70.6--70.8 & 70.7--70.8 \\
Odometry error, max [m]     & 0.106--0.117 & 0.174--0.193 & 0.172--0.198 \\
Thruster peak [N] (of 28.75)& 2.99 & 3.84 & 3.84 \\
Saturation / watchdog trips & 0 / 0 & 0 / 0 & 0 / 0 \\
\bottomrule
\end{tabular}
\end{table}


Three failures from an earlier attempt are retained on record because two of
them are integration lessons that the physical stack will also face: a
middleware session that started isolated and produced a vehicle that never
moved, and a startup race in which the planner engaged before the navigation
estimator was publishing, producing a pose-less maneuver that diverged. Both
are now prevented by a graph-liveness gate and by gating perception on
navigation liveness. The third is a genuine weakness of the committed policy
under an early detection silence, and it is the motivating failure for the
temporal-estimation block rather than something the estimators were credited
with fixing after the fact.

\subsection{Temporal-Estimation Ablation}

The closed-loop ablation ran 100 runs with frozen parameters
(Table~\ref{tab:phase7b}). Two observations dominate. First, the methods are
\emph{indistinguishable on every continuous metric under nominal conditions}:
median clearance spans 0.708--0.713\,m, final lateral error 0.014--0.019\,m,
maneuver time 69.5--69.7\,s, and confirmation delay 1.01--1.03\,s across all
four variants. Second, all of the separation lives in the outlier scenario
E4, where the Kalman family scores 5/5 and 4/5 against 3/5 for both
non-probabilistic variants. The three grazing collisions in E4 occurred at
0.288--0.297\,m minimum clearance against a 0.30\,m threshold, with maneuvers
of normal shape but roughly 0.4\,m smaller lateral amplitude; all three fall
in the hold family and none in the Kalman family. Aggregate success is
therefore 25/25 for T2, 24/25 for T3, and 22/25 for both T0 and T1, and the
single T3 shortfall is an infrastructure freeze rather than a method
difference --- T3 is identical to T2 by construction here.

%==============================================================================
% Phase-7B closed-loop temporal ablation.  Source: docs/PHASE7B_RESULTS.md,
% experiments/simulation/temporal_estimators_phase7b/closedloop/aggregate.json
% (100/100 runs, parameters frozen per D-012, protocol pre-registered).
% Status: READY.  Oracle perception.
%==============================================================================
\begin{table}[t]
\centering
\caption{Closed-loop temporal-estimation ablation: successes out of five
fresh-process runs per cell, under the pre-registered success definition
(no collision, clearance $>0.3$\,m, exactly one avoidance entry, no side
switch, no abort, returned to line). T3 is identical to T2 here ($\theta=0$).
All separation is in E4; E0--E3 are flat.}
\label{tab:phase7b}
\footnotesize
\setlength{\tabcolsep}{4pt}
\renewcommand{\arraystretch}{1.15}
\begin{tabular}{@{}p{34mm} c c c c@{}}
\toprule
\textbf{Scenario} & \textbf{T0} & \textbf{T1} & \textbf{T2} & \textbf{T3} \\
\midrule
E0 clean                        & 5/5 & 5/5 & 5/5 & 5/5 \\
E1 early 2\,s silence           & 5/5 & 4/5 & 5/5 & 5/5 \\
E2 mid-maneuver 2\,s silence    & 4/5 & 5/5 & 5/5 & 5/5 \\
E3 repeated silences            & 5/5 & 5/5 & 5/5 & 5/5 \\
E4 young-track outlier          & 3/5$^{\dagger}$ & 3/5$^{\dagger\dagger}$ & \textbf{5/5} & 4/5 \\
\midrule
\textbf{Total}                  & 22/25 & 22/25 & \textbf{25/25} & 24/25 \\
\bottomrule
\end{tabular}
\vspace{0.5mm}
\raggedright\scriptsize
$^{\dagger}$ one collision. $^{\dagger\dagger}$ two collisions. All three E4
collisions are grazing passes at 0.288--0.297\,m against the 0.30\,m
threshold, all in the hold family, none in the Kalman family. T3's single
shortfall is an infrastructure freeze, not a method difference.
\end{table}


The prediction result runs in the opposite direction, and we report it as a
first-class finding rather than a caveat. During the dropout that begins at
maneuver onset, the simple hold predicts the obstacle's image-space position
substantially better than the constant-velocity filter: prediction-only
centre RMSE is 0.038 versus 0.084 in the early-dropout case and 0.036 versus
0.071 in the mid-maneuver case (Table~\ref{tab:onset},
Fig.~\ref{fig:onset}). The mechanism is measurable in the reference signal:
the obstacle's image-space centre velocity is zero before the vehicle
commits and $-0.032$\,units/s afterwards, so the constant-velocity model
extrapolates a velocity that the maneuver has just invalidated, while the
hold makes no velocity assumption at all. An exploratory damped
constant-velocity variant lands between the two (0.055 and 0.043) and does
not beat the hold; it was discarded as a benchmark method and is recorded
only as a direction for a maneuver-aware model.

%==============================================================================
% Maneuver-onset prediction.  Source:
% experiments/simulation/temporal_estimators_phase7b/onset_analysis/
%   onset_analysis.json  (generated 2026-08-11)
% and experiments/simulation/temporal_estimators/replay_eval/metrics.json.
% Status: READY.
%==============================================================================
\begin{table}[t]
\centering
\caption{Prediction-only centre RMSE (normalized image units) during a
detector dropout, deterministic replay. The hold predicts better than the
constant-velocity filter in both onset cases, because the maneuver
invalidates the velocity the filter extrapolates. T2b is an exploratory
damped-CV variant, reported for completeness and not part of the benchmark.}
\label{tab:onset}
\footnotesize
\setlength{\tabcolsep}{5pt}
\renewcommand{\arraystretch}{1.15}
\begin{tabular}{@{}l c c c@{}}
\toprule
\textbf{Case} & \textbf{T1 (hold)} & \textbf{T2 (CV Kalman)} &
\textbf{T2b (damped)} \\
\midrule
D5 dropout at onset      & \textbf{0.038} & 0.084 & 0.055 \\
D6 mid-maneuver dropout  & \textbf{0.036} & 0.071 & 0.043 \\
\midrule
\multicolumn{4}{@{}l@{}}{\emph{Reference image-space centre velocity:}
$0.000 \rightarrow -0.032$ units/s at commitment} \\
\midrule
\multicolumn{4}{@{}l@{}}{\emph{Outlier case D9 (mature track), overall centre RMSE:}} \\
\quad T0 0.0199 \quad T1 0.0102 & \multicolumn{3}{l}{\quad \textbf{T2 0.0074} $\;\;\Rightarrow\;$ $\chi^2$ gating wins here} \\
\bottomrule
\end{tabular}
\end{table}

%==============================================================================
% Figure --- T1 vs T2 at maneuver onset (case D5).  Status: READY.
% Image regenerated by scripts/analyze_onset.py from
% experiments/simulation/temporal_estimators_phase7b/onset_analysis/.
% Copied to paper/figures/fig_onset_D5.png (do not hand-edit).
%==============================================================================
\begin{figure}[t]
\centering
\includegraphics[width=\columnwidth]{fig_onset_D5.png}
\caption{Image-space obstacle centre during a detector dropout that begins at
maneuver onset (case D5, deterministic replay). Before commitment the
obstacle is stationary in the image; commitment introduces a step in
image-space velocity to $-0.032$\,units/s. The constant-velocity filter
extrapolates the pre-maneuver velocity through the gap and overshoots, while
the hold makes no velocity assumption. Prediction-only RMSE is 0.038 for the
hold against 0.084 for the constant-velocity filter
(Table~\ref{tab:onset}).}
\label{fig:onset}
\end{figure}


Taken together, these two results say something more specific than ``the
Kalman filter helps''. The Kalman machinery earns its place through
$\chi^2$-gated outlier rejection, not through prediction accuracy: it is
outlier armour rather than a better predictor. That distinction matters for
the sim-to-real study, because outlier robustness and prediction accuracy are
stressed by different physical phenomena, and a simulator that reproduces one
faithfully need not reproduce the other.

\subsection{Noise Floor}

The residual failure rate is approximately one non-success per 25 runs from
single-run tail events, plus a roughly 2\% infrastructure mode in which the
vehicle remains pinned near its start. Both are counted as failures under the
pre-registered rules rather than reclassified. This floor is low enough to
compare planners on aggregate success and clearance, and it bounds what the
subsequent blocks can resolve: per-scenario differences smaller than about
two runs in five are not resolvable at this sample size, which is the reason
the planner campaign uses ten repetitions per cell.
